% GENERATED FILE. Edit canonical lessons, then run npm run generate:books.
% canonical-source-hash: fnv1a64:d6dd5fb63c899777
% canonical-lessons: KA-C47-leg, KA-C47-tooth, KA-C47-hair, KA-C47-finger, KA-C47-stomach

\chapter{The Leg and the Tooth}
\label{ch:ka-body-more}
\hlchaptermodality{47}

\begin{chapteropening}
\textbf{By the end of this chapter:} \emph{I can name five more parts of the body in Kannada, and point to the two regular sound exchanges that separate them from their Tamil and Telugu cousins.}

Complete KA-C47-stomach and say all five words in order.
\end{chapteropening}

\section[kālu]{\kn{ಕಾಲು} (kālu) — leg, foot}

\label{lesson:KA-C47-leg}



\begin{quote}
Before the new run: what did \emph{pustaka} mean, and what did \emph{uppu} mean?
\end{quote}

\subsection*{You'll want to know: \kn{ಕಾಲು}}
\textbf{\kn{ಕಾಲು}} (\emph{kālu}) — \textquotedblleft{}leg, foot\textquotedblright{}.

Proto-Dravidian \textbf{*kāl}, and every sister kept it nearly whole: Tamil \ta{கால்}, Telugu \te{కాలు}, Malayalam \ml{കാൽ}.

Set that beside \kn{ಕೈ}, the hand-word, which wore down so far that the family resemblance is hard to see at all. Two body words of the same age; one held its shape and one did not, and nothing about the meaning predicts which.

\kn{ಕಾಲು} covers the whole limb, foot and leg together. It also names a quarter — a fourth of anything — so \kn{ಕಾಲು} \kn{ಗಂಟೆ} is a quarter of an hour. Nothing joins the leg to the fraction but an accident of sound, and context is the only fix, which is what context is for.

\begin{grammarlens}[title={What you've built}]
The first of five more body words.
\end{grammarlens}

\subsection*{Your turn}
Say these aloud:

\begin{itemize}
\raggedright
  \item \emph{kālu}
  \item it once more, slowly
  \item \emph{kālu}, then \emph{tale}, so the top and the bottom sit together
\end{itemize}

\begin{tcolorbox}[breakable,colback=teal!4,colframe=teal!35!black,title={Before you move on}]
What does \kn{ಕಾಲು} mean? (\textquotedblleft{}leg, foot\textquotedblright{}.) Say it once more.
\end{tcolorbox}
\section[hallu]{\kn{ಹಲ್ಲು} (hallu) — a tooth}

\label{lesson:KA-C47-tooth}



\begin{quote}
Before the new one: what did \emph{kālu} mean?
\end{quote}

\subsection*{You'll want to know: \kn{ಹಲ್ಲು}}
\textbf{\kn{ಹಲ್ಲು}} (\emph{hallu}) — \textquotedblleft{}a tooth\textquotedblright{}.

Proto-Dravidian \textbf{*pal}. Tamil kept \ta{பல்}; Telugu took a third road and came out with \te{పన్ను}; Kannada applied its own law and got \kn{ಹಲ್ಲು}, the \emph{p-} turned to \emph{h-} exactly where the chapter before said to watch for it.

That makes \kn{ಹಣ್ಣು} and \kn{ಹಲ್ಲು} a matched pair: two inherited words, one \emph{h} apiece, and a \emph{p} at the front of both a thousand years ago.

\begin{grammarlens}[title={What you've built}]
Two: a leg and a tooth.
\end{grammarlens}

\subsection*{Your turn}
Say these aloud:

\begin{itemize}
\raggedright
  \item \emph{hallu}
  \item it once more, slowly
  \item \emph{kālu}, then \emph{hallu}, and say which of the two also names a fraction
\end{itemize}

\begin{tcolorbox}[breakable,colback=teal!4,colframe=teal!35!black,title={Before you move on}]
What does \kn{ಹಲ್ಲು} mean? (\textquotedblleft{}a tooth\textquotedblright{}.) Say it once more.
\end{tcolorbox}
\section[kūdalu]{\kn{ಕೂದಲು} (kūdalu) — hair}

\label{lesson:KA-C47-hair}



\begin{quote}
Before the new one: what did \emph{hallu} mean?
\end{quote}

\subsection*{You'll want to know: \kn{ಕೂದಲು}}
\textbf{\kn{ಕೂದಲು}} (\emph{kūdalu}) — \textquotedblleft{}hair\textquotedblright{}.

\kn{ಕೂದಲು} is the hair on a head, counted as a mass rather than strand by strand. Tamil's \ta{கூந்தல்} (\emph{kūntal}) is the same inherited word.

Look at what it finishes on. \kn{ಕಾಲು}, \kn{ಹಲ್ಲು}, \kn{ಕಣ್ಣು} and \kn{ಹಣ್ಣು} all end on a small \emph{-u} as well. Kannada dislikes leaving a word standing on a bare consonant and sets a light vowel behind it — a prop, carrying no meaning of its own. Tamil tolerates the bare consonant and writes \ta{பல்} and \ta{கால்} with nothing after them at all.

Careful Kannada also has \kn{ಕೇಶ} (\emph{kēśa}), straight from Sanskrit, for the barber's signboard and the poem.

\begin{grammarlens}[title={What you've built}]
Three: a leg, a tooth, and hair.
\end{grammarlens}

\subsection*{Your turn}
Say these aloud:

\begin{itemize}
\raggedright
  \item \emph{kūdalu}
  \item it once more, slowly
  \item all three, and listen for the little \emph{-u} at the end of each
\end{itemize}

\begin{tcolorbox}[breakable,colback=teal!4,colframe=teal!35!black,title={Before you move on}]
What does \kn{ಕೂದಲು} mean? (\textquotedblleft{}hair\textquotedblright{}.) Say it once more.
\end{tcolorbox}
\section[beraḷu]{\kn{ಬೆರಳು} (beraḷu) — a finger}

\label{lesson:KA-C47-finger}



\begin{quote}
Before the new one: what did \emph{kūdalu} mean?
\end{quote}

\subsection*{You'll want to know: \kn{ಬೆರಳು}}
\textbf{\kn{ಬೆರಳು}} (\emph{beraḷu}) — \textquotedblleft{}a finger\textquotedblright{}.

Tamil has \ta{விரல்} (\emph{viral}), Malayalam \ml{വിരൽ}, Telugu \te{వేలు} (\emph{vēlu}) — and Kannada opens the word with \emph{b-} where its sisters have \emph{v-}.

That is the second regular exchange this book has shown you, the one already visible in \kn{ಬರೆ}, 'to write', where the sisters again have a \emph{v}. One family, two swaps: \emph{p-} for \emph{h-}, \emph{v-} for \emph{b-}. Between them they account for a great many words that look unrelated at first glance and are nothing of the kind.

\begin{grammarlens}[title={What you've built}]
Four. A leg, a tooth, hair, a finger.
\end{grammarlens}

\subsection*{Your turn}
Say these aloud:

\begin{itemize}
\raggedright
  \item \emph{beraḷu}
  \item it once more, slowly
  \item \emph{beraḷu}, then \emph{bare}, and hear the same \emph{b} at the front of both
\end{itemize}

\begin{tcolorbox}[breakable,colback=teal!4,colframe=teal!35!black,title={Before you move on}]
What does \kn{ಬೆರಳು} mean? (\textquotedblleft{}a finger\textquotedblright{}.) Say it once more.
\end{tcolorbox}
\section[hoṭṭe]{\kn{ಹೊಟ್ಟೆ} (hoṭṭe) — the stomach, the belly}

\label{lesson:KA-C47-stomach}



\begin{quote}
Before the new one: what did \emph{beraḷu} mean?
\end{quote}

\subsection*{You'll want to know: \kn{ಹೊಟ್ಟೆ}}
\textbf{\kn{ಹೊಟ್ಟೆ}} (\emph{hoṭṭe}) — \textquotedblleft{}the stomach, the belly\textquotedblright{}.

Telugu has \te{పొట్ట} (\emph{poṭṭa}) for the same thing, the \emph{p} being the older shape and Kannada's \emph{h} the newer one. Tamil went its own way here and reaches for \ta{வயிறு} (\emph{vayiṟu}), a word from another root entirely — a reminder that the sisters agree often, not invariably.

Five body words, and three of them — \kn{ಹಲ್ಲು}, \kn{ಹೊಟ್ಟೆ}, and the fruit from the chapter before — carry the same \emph{h} into the same first slot.

\begin{grammarlens}[title={What you've built}]
Five more body words: a leg, a tooth, hair, a finger, the belly.
\end{grammarlens}

\subsection*{Your turn}
Say these aloud:

\begin{itemize}
\raggedright
  \item \emph{hoṭṭe}
  \item it once more, slowly
  \item all five in order, then the three that begin with \emph{h}
\end{itemize}

\begin{tcolorbox}[breakable,colback=teal!4,colframe=teal!35!black,title={Before you move on}]
What does \kn{ಹೊಟ್ಟೆ} mean? (\textquotedblleft{}the stomach, the belly\textquotedblright{}.) Say it once more.
\end{tcolorbox}
